Qt --- это кросс-платформенная библиотека классов C++ для создания графических пользовательских интерфейсов (GUI) компании Digia. Эта библиотека полностью объектно-ориентированная, что обеспечивает лёгкое расширение возможностей и создание новых компонентов. Кроме того, она поддерживает большое количество платформ.

Qt позволяет запускать написанное с его помощью ПО в большинстве современных операционных систем путём простой компиляции программы для каждой ОС без изменения исходного кода. Включает в себя все основные классы, которые могут потребоваться при разработке прикладного программного обеспечения, начиная от элементов графического интерфейса и заканчивая классами для работы с изображениями, видео и мультипоточными приложениями. Qt является полностью объектно-ориентированным, легко расширяемым и поддерживающим технику компонентного программирования.

Список использованных классов фреймворка Qt:
\begin{itemize}
\item QDebug;
\item QDir;
\item QFile;
\item QTextStream;
\item QString;
\item QImage;
\item QTime.
\end{itemize}

Класс QDebug обеспечивает выходной поток для отладочной информации.

Класс QDir обеспечивает доступ к структуре каталогов и их содержимого.

Класс QFile предоставляет интерфейс для чтения и записи файлов.

Класс QTextStream предоставляет удобный интерфейс для чтения и записи текста.

Класс QString обеспечивает строку символов Unicode.

Класс QImage предоставляет аппаратно-независимую работу с изображениями, даёт прямой доступ к данным каждого пикселя и может быть использован в качестве устройства рисования. На ранних стадиях создания ПО этот класс использовался в качестве основного инструмента для работы с изображениями, но впоследствии при переходе на основные модули библиотеки DV стал использоваться объект Data2D и методы работы с ним.

Класс QTime обеспечивает функции таймера. В проекте используется для оценки быстродействия разрабатываемого модуля при разных входных параметрах~\cite{qt}.

Для совместимости с библиотекой DV решено использовать систему сборки CMake взамен QMake. CMake широко используется с Qt.
Список классов и методов библиотеки DV, используемых в проекте:
\begin{itemize}
\item Data2D;
\item Matx22d;
\item Vec2d;
\item VF2d.
\end{itemize}

Список стандартных библиотек используемых в проекте:

\begin{itemize}
\item getopt.h;
\item iostream;
\item math.h;
\item stdlib.h;
\item vector.
\end{itemize}

getopt библиотечная функция, специально разработанная для того чтобы облегчить обработку входных команд.

iostream заголовочный файл с классами, функциями и переменными для организации ввода-вывода в языке программирования C++. Он включён в стандартную библиотеку C++. Название образовано от Input/Output Stream (``поток ввода-вывода'').

math.h — заголовочный файл стандартной библиотеки языка программирования С, разработанный для выполнения простых математических операций.

stdlib.h — заголовочный файл стандартной библиотеки языка Си, который содержит в себе функции, занимающиеся выделением памяти, контролем процесса выполнения программы, преобразованиями типов и др.

vector — это замена стандартному динамическому массиву, память для которого выделяется вручную с помощью оператора new.
